% P10-HITL, 15.09.2026: Entscheidungszweck und belegte Zustandswirkung präzisiert.
% P10-4b, 14.09.2026: Begriffs-/Beleganschlüsse und abgegrenzte Prüfaussagen präzisiert.
% M19 (12.09.2026): Leserführung und präzise Ergebnisdarstellung;
% Messwerte und historische Nachweisgrenzen bleiben erhalten.
% ============================================================
% §7.4 Kontrollen der Wahrheits-Mutation — v01 (Claude, 07.09.)
% NEUE 8er-Struktur (Umsetzungsplan §5: 5.1 Bestandsevidenz nach
% Assertions [Z3] · 5.2 T1–T4 am realen Apply-Pfad · 5.3 Maker–
% Checker-Nullbefund). Budget ~1,5 S. + 2 Tabellen.
% Leitplanken: Assertions statt Testnamen · Zustandswirkung als
% Beleg · G04 offen benennen · keine Erfolgsquote/Bestehensschwelle.
% ============================================================
% HERKUNFT: Z3-Zuordnung = thesis-evidence/w2/z3-assertion-zuordnung.md
% (36 Methoden, Prüftiefen 28/5/2/1; inkl. Korrektur G01/G02 nach
% Kollegen-Review 07.09.); Matrix = runs/e2e-evidenz/
% fehlerklassen-matrix.md (v2, Bilanz 6/4/2); T1–T4 =
% thesis-evidence/w2/t1-t4-kontrolltests.md + Testcode
% AgenticSdlc.Tests/FullWorkflow/IngestionApplyGovernanceTests.cs
% (Beleg-Lauf 07.09.: Filter 4/4, Suite 688/688); Signal-Herkunft =
% z1-stufenanalyse.md §4 (Lauf 20260905_164650_8189ea);
% canonical-gate = runs/ledger/20260905_164650_8189ea/gate/.
% ------------------------------------------------------------

\section{Kontrollen von Änderungen am Core}
\label{sec:eval-kontrollen}

Im vorgesehenen Anwendungsablauf werden Änderungen am fachlichen
Projektbestand über die gemeinsame Speichernaht mit strukturellen
Invarianten gespeichert. Die menschliche
Einzelfreigabe fachlicher Übernahmen gilt für die
Fortschreibungspfade des vorhandenen Bestands --- der initiale Seed
der Bestandsherstellung ist davon ausgenommen
(Abschnitt~\ref{sec:endtoend-pfade}). Das umfassende
Autorisierungsziel A6 ist somit nur teilweise umgesetzt; dies
begrenzt auch die Autorisierungsannahme des Projektzustands in A5,
A7 und A8. Prüfgegenstand ist die Bindung der fachlichen
Zustandsänderung an die vorgesehene Freigabe: Eine nicht freigegebene
Operation darf nicht übernommen werden; bei einem gültigen und
freigegebenen Vorschlag muss die vorgesehene Übernahme möglich sein.
Die Untersuchung bewertet
damit die Umsetzung der Entscheidungshoheit, nicht die fachliche
Richtigkeit menschlicher Urteile. Der Abschnitt prüft die
realisierten Kontrollen auf zwei Wegen: durch die
Zuordnung der \emph{bestehenden} Tests deterministischer Kontrollen
anhand ihrer tatsächlich geprüften Bedingungen (Assertions) und durch
vier \emph{neu ausgeführte} Kontrolltests am realen
Anwendungspfad.

\subsection{Bestehende Tests nach Prüfinhalt}
\label{sec:eval-z3}

36 vorhandene Testmethoden decken die Prüfmechanik der Kette
(Kanonisierungs-Check, Referenzprüfung und Rückstufung zur
Nachprüfung), die Invariantenprüfung am Core, die Snapshot-Historie
und die deterministischen Projektionen ab. Die Realisierung der
Invariantenprüfung als \texttt{CoreKangal} erläutert
Abschnitt~\ref{sec:realisierung-persistenz-aussenkopplung}. Die Zuordnung der Testmethoden erfolgte anhand
der tatsächlichen Prüfinhalte; dabei wurde je Methode die
Prüftiefe erfasst: 28 Methoden prüfen Rückgabewerte im Arbeitsspeicher (In-Memory), fünf
den echten Datei-/Repository-Zustand, zwei einen Ausnahme-Wurf und
eine kombiniert Ausnahme und Dateizustand. Diese Kartierung
korrigierte zwei frühere Zuordnungsfehler: Ein Test prüft entgegen
seinem Namen nur die zurückgegebene Dateiliste, keinen
Datenträger-Schreibvorgang; und die Challenge „unbekannte
Relation`` war fälschlich als blockiert-getestet geführt ---
tatsächlich quittiert die Invarianten-Prüfung unbekannte
Relationstypen nur mit einer protokollierten Warnung, und kein Test prüft
dieses Verhalten. Der einzige Test, der die strukturelle Blockade
der Speichernaht am echten Dateizustand nachweist, ist der
Save-Test der Invarianten-Prüfung (wirft und schreibt nichts,
anschließend persistiert ein gültiger Save).

\subsection{Vier Kontrolltests am realen Anwendungspfad}
\label{sec:eval-t1t4}

Die von der Testkartierung offengelegten Lücken --- Apply ohne
Freigabe und die Zustandswirkung einer Ablehnung waren nirgends am
persistierenden Pfad getestet --- wurden durch vier neue Tests
geschlossen (Tabelle~\ref{t:eval-kontrolltests}). Alle vier
arbeiten ohne simulierte Ersatzkomponenten (Mocks): Sie setzen einen realen Core in einem
Temporärverzeichnis auf, durchlaufen die geteilte
Anwendungsroutine des Ingestion-Tors
(Apply; gemeinsam von CLI und Workflow genutzt) und lesen den Core-Zustand anschließend aus
dem Dateisystem zurück. Unterschieden wird dabei der fachliche
Item-Bestand von den gespeicherten Ablehnungsvermerken: Eine abgelehnte
Änderung darf keinen fachlichen Inhalt erzeugen, wohl aber einen
begründeten Ablehnungsvermerk hinterlassen.

\input{\tabledir/eval-kontrolltests.tex}

In den geprüften Fällen verhinderten fehlende beziehungsweise
abgelehnte Freigaben die fachliche Übernahme; die erneute
Anwendung desselben Plans erzeugte keine zweite fachliche Wirkung;
der gültige Kontrollfall bestätigte die vorgesehene Übernahme
einschließlich Herkunft und Historisierung --- und belegt damit
zugleich, dass die negativen Fälle nicht deshalb folgenlos
blieben, weil der Pfad grundsätzlich nichts speichert. Die
Aussagen gelten für den geprüften Anwendungspfad; sie sind ein
anderer Nachweis als die strukturelle Naht-Blockade und ersetzen
weder diese noch eine allgemeine Absicherung gegen doppelte
externe Schreibvorgänge nach einem Absturz.

\subsection{Bilanz und Grenzen}
\label{sec:eval-kontrollbilanz}

Tabelle~\ref{t:eval-challenge} schließt den zwölfteiligen
Challenge-Katalog aus der Anforderungslandkarte: sechs Fälle sind
getestet, vier teilweise belegt, zwei nicht getestet. Offen
bleiben die unbekannte Relation (nur Warnung, ungetestet) und der
direkte Umgehungsweg außerhalb der Speichernaht (G04) --- Aussagen
wie „technisch nicht umgehbar`` sind daher auf das belegte
Architektur- und Ausführungsmodell begrenzt (eine Speichernaht,
neun inventarisierte Schreiber). Die geprüften Governance-Pfade
belegen zudem keine durchgängige menschliche Einzelautorisierung
sämtlicher Bestandsinhalte: Der initiale Seed der
Bestandsherstellung schreibt ohne Einzelautorisierung der
entstehenden Elemente (Abschnitt~\ref{sec:endtoend-pfade}); die
Kontrollnachweise betreffen die Fortschreibungspfade. Eine
Erfolgsquote oder
Bestehensschwelle wird bewusst nicht gebildet.

\input{\tabledir/eval-challenge-bilanz.tex}

Für die Prüf-Reparatur-Schleife der Kanonisierung gilt: Im
ausgewerteten Stresslauf passierte der kanonische Bestand den
deterministischen Übergangs-Check im ersten Versuch; dieser
Reparaturzweig wurde nicht ausgelöst, sein Korrektureffekt ist
daher nicht gemessen --- ein nicht ausgeführter Mechanismus ist
kein negativer Wirksamkeitstest. Davon zu unterscheiden sind zwei
weitere Mechanismen desselben Laufs, die regulär ausgeführt
wurden: Die Facetten-Validierung meldete 22 Facettenbefunde
(20 mit angewandter Facetten-Reparatur) und stufte 18 Claims zur
Nachprüfung ein; ein Referenz-Reparatur-Durchgang der
Hinweisstufe lief einmal (Stufenzuordnung:
Anhang~\ref{anh:evaluation}). Deren fachliche Wirkung wurde nicht
separat bewertet.

Der deterministische Übergangs-Check der Kanonisierung ist
zugleich die in Abschnitt~\ref{sec:eval-stufen} belegte Grenze:
Er prüft referenzielle Deckung, nicht semantische Erhaltung ---
die beiden dort nachgewiesenen Merge-Verluste passierten ihn
unbeanstandet. Die Erkennung definierter Fehlerklassen durch die
Schleife ist deterministisch getestet (erste Tabellengruppe der
Kartierung). Ein Korrektureffekt der LLM-Reparatur wurde in der
hier ausgewerteten Stufenanalyse nicht gemessen. Der historische
Backlog-Lauf \texttt{20260804\_074419\_acab74} zeigt dagegen einen
Reparaturversuch nach zunächst ausgebliebener Gesamteinreichung.
Innerhalb dieses Versuchs wurden vier fehlende Verbindungen eines
19-PBI-Entwurfs ergänzt; anschließend bestand die Strukturprüfung
(Abschnitt~\ref{sec:saeule-semantische-transformation}, E45-02).
Dieser formative Einzelfall belegt eine strukturelle Korrektur,
keine semantisch geprüfte Reparaturqualität oder allgemeine
Wirksamkeit des Musters.
